\section{System Model and Threat Model}
\label{sec:system}

\subsection{Signal Representation and Dataset}

We consider complex baseband IQ signals represented as real-valued tensors
$\mathbf{x} \in \mathbb{R}^{2 \times T}$, where the first dimension indexes
the in-phase (I) and quadrature (Q) channels and $T$ is the number of time
samples. Each channel captures a different aspect of the baseband signal; their
joint representation encodes the amplitude and phase of the modulated carrier.
Following the RML2016.10a dataset convention~\cite{otoole2016rml}, we use
$T = 128$ samples, corresponding to 16 symbols at 8 samples-per-symbol with
a root-raised-cosine pulse shape.

The RML2016.10a dataset contains 220,000 IQ bursts drawn from 11 modulation
classes: BPSK, QPSK, 8PSK, QAM16, QAM64, PAM4, AM-DSB, AM-SSB, WBFM, GFSK,
and CPFSK. Each burst is labeled with the modulation type and the signal-to-noise
ratio (SNR) at which it was recorded, spanning $-20$ to $+18$~dB in 2~dB steps.
At high SNR ($\geq 0$~dB), the AWN classifier~\cite{li2023awn} achieves above
91\% overall accuracy across all 11 classes.

\subsection{AWN Classifier Architecture}

The Adaptive Wavelet Network (AWN)~\cite{li2023awn} takes batched IQ tensors of
shape $[N, 2, T]$ and predicts one of $C = 11$ modulation classes. The forward
pass has three stages:

\textbf{Convolutional feature extraction.} A 2D convolution with kernel size
$[2, 7]$ integrates the I and Q channels, yielding a feature map of shape
$[N, 64, 1, T]$, which is then squeezed to $[N, 64, T]$. A second 1D
convolution with batch normalization~\cite{ioffe2015batch} and LeakyReLU
extracts temporal features, producing $[N, 64, T]$ feature maps.

\textbf{Adaptive wavelet decomposition.} The core module applies a
learnable lifting scheme~\cite{li2023awn} that recursively splits the
feature map into even and odd subsequences, applies learnable Predictor ($P$)
and Updator ($U$) operators to produce approximation coefficients $\mathbf{c}$
and detail coefficients $\mathbf{d}$, and repeats for \texttt{num\_levels}
levels. The detail regularization loss is:
\begin{equation}
  \mathcal{R}_d = \sum_{\ell=1}^{L} \lambda_d \,\|\mathbf{d}_\ell\|_2^2,
  \quad
  \mathcal{R}_c = \lambda_c \,|\bar{\mathbf{c}}|,
\end{equation}
where $\bar{\mathbf{c}}$ is the mean of the final approximation. The total
training loss is:
\begin{equation}
  \mathcal{L} = \mathcal{L}_{\mathrm{CE}} + \mathcal{R}_d + \mathcal{R}_c,
\end{equation}
where $\mathcal{L}_{\mathrm{CE}}$ is the cross-entropy classification loss.

\textbf{Attention and classification.} Squeeze-excitation (SE)
attention~\cite{hu2018squeeze} computes per-level weights for the multi-scale
approximation and detail coefficient tensors. The attention-weighted features
are pooled and passed through two fully-connected layers with dropout and
LeakyReLU activations to produce the final class logits. The model is trained
using the Adam optimizer~\cite{kingma2015adam} with early stopping.

\subsection{Attack Models}

We evaluate four attacks that span the gradient-based and optimization-based
attack families. All attacks use \texttt{torchattacks}~\cite{kim2020torchattacks}
with per-sample minmax normalization: each IQ sample is mapped to $[0, 1]$ using
$\tilde{x} = (x - x_{\min}) / (x_{\max} - x_{\min})$ before attack generation
and mapped back afterward. This makes epsilon values interpretable relative to
the signal range.

\textbf{FGSM.} The Fast Gradient Sign Method~\cite{goodfellow2015fgsm}
computes a single gradient step:
\begin{equation}
  \mathbf{x}_{\mathrm{adv}} = \mathbf{x} + \varepsilon \cdot
    \mathrm{sign}\bigl(\nabla_{\mathbf{x}}\,\mathcal{L}_{\mathrm{CE}}
    (\mathbf{x}, y)\bigr),
\end{equation}
where $\varepsilon$ bounds the $L_\infty$ perturbation magnitude. FGSM
is computationally efficient but produces suboptimal adversarial examples.

\textbf{PGD.} Projected Gradient Descent~\cite{madry2018pgd} iterates
multiple gradient steps with projection onto the $\varepsilon$-ball:
\begin{equation}
  \mathbf{x}^{(t+1)} = \Pi_{\mathcal{B}_\varepsilon(\mathbf{x})}
    \!\left(\mathbf{x}^{(t)} + \alpha\,\mathrm{sign}\bigl(
    \nabla_{\mathbf{x}}\,\mathcal{L}_{\mathrm{CE}}(\mathbf{x}^{(t)}, y)
    \bigr)\right),
\end{equation}
where $\Pi$ denotes projection and $\alpha$ is the step size. PGD with a
random start is regarded as a strong first-order attack~\cite{croce2020reliable}.

\textbf{CW ($L_2$).} The Carlini--Wagner attack~\cite{carlini2017towards}
solves the constrained optimization:
\begin{equation}
  \min_{\boldsymbol{\delta}}\; \|\boldsymbol{\delta}\|_2 + c\, f(\mathbf{x} + \boldsymbol{\delta}),
\end{equation}
where $f(\cdot)$ is a surrogate misclassification objective and $c$ balances
perturbation norm against attack success. CW distributes adversarial energy
broadly across spectral frequencies, making it resistant to simple band-limited
filtering. We use $c = 10$, 200 optimization steps.

\textbf{EAD ($L_1$ and EN).} The Elastic-Net Attack on Deep Neural
Networks~\cite{chen2020ead} adds an $L_1$ regularization term to the CW
objective:
\begin{equation}
  \min_{\boldsymbol{\delta}}\; \beta\|\boldsymbol{\delta}\|_1 +
    \|\boldsymbol{\delta}\|_2 + c\, f(\mathbf{x} + \boldsymbol{\delta}),
\end{equation}
where $\beta$ controls the sparsity of the perturbation. We evaluate both the
$L_1$ variant (EADL1, $\beta$ large) and the elastic-net variant (EADEN,
$\beta$ moderate). EADL1 produces sparse perturbations concentrated in few
spectral components, while EADEN produces a mixture of sparse and distributed
perturbations.

\subsection{Threat Model}

\textbf{Attacker capabilities.} We consider a white-box adversary with full
knowledge of the AWN classifier architecture, trained weights, and input
normalization procedure. The attacker can query the classifier and compute
exact gradients. This is the strongest adversary assumption and establishes
a worst-case defense baseline.

\textbf{Attacker goal.} The adversary seeks untargeted misclassification: any
prediction other than the true modulation class is considered a success. The
attacker does not attempt to cause the classifier to predict any specific wrong
class.

\textbf{Attacker constraint.} Perturbations are bounded by the chosen
attack norm: $\|\boldsymbol{\delta}\|_\infty \leq \varepsilon$ for FGSM/PGD,
$\|\boldsymbol{\delta}\|_2 \leq$ (implicitly, via CW minimization) for CW, and
elastic-net bounded for EAD. We evaluate $\varepsilon \in \{0.05, 0.1\}$ in
the per-sample minmax-normalized space (equivalent to 5--10\% of the
per-sample signal range).

\textbf{Defense-unaware model.} We assume the attacker does not adapt to the
defense mechanism: adversarial examples are generated against the undefended
AWN classifier. This models the realistic deployment scenario where the defense
is applied transparently at the receiver and is not publicly disclosed. Adaptive
attacks against the defense are left for future work.

\textbf{Defender capabilities.} The defender can preprocess each IQ burst
before classification, with access only to the received signal. The defender
does not know the true modulation class, the attack type, or the attack
parameters. The defense must operate in real time: latency must be sub-millisecond
per sample to meet the requirements of practical AMC deployment contexts,
including ITU/FCC spectrum monitoring stations, CBRS Environmental Sensing
Capability (ESC) nodes, and military electronic support measures (ESM)
systems~\cite{proakis2008digital}. In each of these contexts, AMC drives
downstream decisions (spectrum allocation, threat identification) that require
near-real-time throughput, typically at 25--100~Msps IQ digitization rates with
burst processing latency budgets of 1--10~ms per channel.
